\chapter{Levels of Rigor: A Methodological Framework}
\label{ch:levels-of-rigor}

\epigraph{\itshape The map is not the territory; but a usable map is
not a territory either.}{}

\noindent
\xref{ch:two-failures} ended with the proposal that the metaphysical
claims of the Islamic tradition deserve theoretical treatment: that
they should be stated precisely, formalized where possible, derived
from first principles where the formalism permits, and tested
empirically where the data allows. This chapter takes up the obvious
next question. Not all claims admit the same kind of treatment. Some
claims can be made precise and derived from formal premises; others
can be expressed in formal vocabulary but not derived; still others
resist formalization but can be honestly characterized as claims
about features of the world. A methodology that fails to distinguish
these cases will either over-claim (asserting causal mechanism where
only formal analogy is available) or under-claim (refusing to
engage with claims because they cannot be reduced to causal mechanism).

This chapter develops the three-level taxonomy that organizes every
substantive treatment in the rest of the book. \xref{sec:why-taxonomy}
explains why a taxonomy is needed at all. \xref{sec:level-1} defines
Level 1 (metaphor); \xref{sec:level-2} defines Level 2 (formal
analogy); \xref{sec:level-3} defines Level 3 (causal mechanism).
\xref{sec:worked-examples} works through five examples that calibrate
the reader's intuitions for the classification. \xref{sec:falsifiability}
discusses falsifiability and what it requires at each level.
\xref{sec:akbarian-role} addresses the special status of Akbarian
metaphysics in the framework. \xref{sec:ahead-2} previews how the
levels apply to the eight canonical claims that Part III will treat.

\section{Why a taxonomy is needed}
\label{sec:why-taxonomy}

The two failures discussed in \xref{ch:two-failures} have a common
diagnosis. The mainstream Islamic economics literature lacks the
methodological apparatus to engage with metaphysical claims. The
literature has, broadly, two registers. The first is the formal
register of mainstream economics: agents have preferences, technologies
are described by production functions, equilibria are characterized
by no-arbitrage conditions. Within this register, metaphysical claims
about provision, blessing, intention, and decree have no place; they
cannot be expressed in the formal vocabulary. The second is the
homiletic register of devotional literature: claims are asserted as
moral truths and the believer is invited to accept them. Within this
register, metaphysical claims have everywhere to be, but the
register provides no resources for testing, formalizing, or
analyzing them.

The two registers are mutually exclusive in practice. An author who
attempts to combine them produces something that is neither
formal economics nor devotional literature: it is the genre of
``theological economics'' that, when done badly, produces concordism
and pop-spirituality, and that, when done well, produces work that
mainstream economists do not engage with and devotional readers do
not find satisfying.

What is missing is a third register: a methodology that permits
\emph{graded engagement} with metaphysical claims. Some claims can
be formalized completely; others partially; others not at all. The
methodology must distinguish these cases, specify what can be done
in each, and prevent the over-claims and under-claims that
characterize the literature when the distinctions are not made.

The three-level taxonomy is this methodology. It is not a doctrine.
It is a procedure. It tells us, for any given claim, what we can do
with it and what we cannot. It tells us which formal apparatus
applies, which empirical predictions follow, and which questions
remain genuinely open. It is the device that converts the
metaphysical content of the tradition from an unprocessed corpus of
homiletic assertion into a structured research program.

\subsection{Antecedents in other fields}

The taxonomy is not original to this work in its broad outline.
Analogous distinctions are routinely made in the philosophy of
science. The distinction between \emph{causal} and \emph{constitutive}
explanation is widely discussed.\footnote{See, classically, Salmon, W.
(1984), \emph{Scientific Explanation and the Causal Structure of the
World}, Princeton University Press; and more recently Craver, C.
(2007), \emph{Explaining the Brain}, Oxford University Press.} The
distinction between \emph{first-order} and \emph{higher-order}
empirical claims is fundamental to the philosophy of physics. The
distinction between \emph{formal} and \emph{material} truth is at
least as old as Aristotle. What is novel here is the application of
these distinctions to the specific problem of articulating Islamic
metaphysical claims in modern theoretical vocabulary, and the
specific three-fold structure that emerges from the application.

Readers familiar with these antecedent distinctions will recognize
the structure of the taxonomy. Readers unfamiliar with them will
find the taxonomy self-contained: every term is defined and every
claim is illustrated.

\section{Level 1: Metaphor}
\label{sec:level-1}

\begin{definition}[Level 1: Metaphor]
\label{def:metaphor}
A metaphysical claim is treated at \emph{Level 1} when it is asserted
as a figurative or pedagogical illustration without specification of
its formal content, its empirical predictions, or its derivability
from premises. A Level 1 treatment may be vivid, memorable, or
spiritually nourishing, but it does not produce theoretical content
that can be tested, derived, or extended.
\end{definition}

The simplest example: the claim ``rizq is like air --- it surrounds
you whether you reach for it or not.'' This is a Level 1 treatment.
The simile is suggestive: it captures something of the believer's
felt sense that provision arrives independently of striving. But
nothing in the simile commits us to a specific empirical claim.
Air is a substance; rizq might or might not be. Air is uniformly
distributed in the atmosphere; rizq might be variably distributed.
Air can be measured by pressure and density; nothing in the simile
tells us how rizq is measured. The simile illustrates an attitude
toward provision; it does not specify the structure of provision.

Level 1 treatments are the natural register of homiletic literature.
The Friday sermon is not the place for derivation. The believer who
hears that ``the gates of provision open for the one who is
grateful'' is being given a moral and spiritual orientation, not a
theoretical claim. Level 1 treatments are appropriate, even necessary,
in their proper context. The error is to confuse them with
theoretical claims, or to suppose that they are the only possible
treatment of the underlying material.

\subsection{What Level 1 commits us to}

A Level 1 treatment commits us to almost nothing in theoretical
terms. It does not commit us to the existence of any entity beyond
those of common sense. It does not commit us to the truth of any
empirical proposition. It does not commit us to the validity of
any inference rule beyond ordinary discursive reasoning. The
treatment functions as illustration; the illustration's success is
measured by its pedagogical effectiveness, not by its theoretical
content.

This is a virtue when the goal is pedagogical and a limitation when
the goal is theoretical. The believer instructed by the simile of
air can take the simile as a guide to attitude. The economist who
wants to test whether provision behaves like a conserved quantity
gets nothing from the simile. The two purposes are different and
the simile serves the first.

\subsection{Level 1 in this book}

The book uses Level 1 treatments rarely, and always with explicit
flagging. When a Level 1 expression is used --- typically in the
introduction of a section, to motivate a more rigorous treatment that
follows --- the use is signaled by phrases like ``the tradition has
sometimes expressed this as'' or ``the homiletic register puts it
this way.'' The reader is alerted that what follows is not the
substantive treatment but the motivating framing. The substantive
treatment will come at Level 2 or Level 3.

\section{Level 2: Formal analogy}
\label{sec:level-2}

\begin{definition}[Level 2: Formal analogy]
\label{def:formal-analogy}
A metaphysical claim is treated at \emph{Level 2} when it is given
mathematical or formal structure that is structurally identical to
some known formal system, without claiming that the known system is
the actual mechanism. The treatment specifies what would be true of
the claim under the formal structure, what relations between
variables it implies, and what consequences it produces, but it does
not assert that the underlying physical or economic mechanism is the
one operating in the analogue system.
\end{definition}

Level 2 is the level at which structural insight becomes possible
without empirical commitment. The treatment imports the formal
structure of a well-understood system to organize the analysis of a
less well-understood phenomenon. The phenomenon may or may not, in
fact, be governed by the same mechanism that governs the analogue
system. What the formal analogy gives us is a rigorous characterization
of \emph{what would be true of the phenomenon if it had the same
structure}.

A canonical example from outside Islamic economics: the use of
field-theoretic vocabulary to describe phenomena in social science.
The ``social field'' of Pierre Bourdieu, the ``electromagnetic field''
of social attraction in popular sociology --- these are formal
analogies. Bourdieu does not claim that social attraction is governed
by Maxwell's equations. He claims that the formal structure of a
field --- a domain in which forces act on positioned actors --- is
useful for organizing the analysis of social phenomena. The analogy
is doing real work: it suggests propositions, organizes the data,
and yields hypotheses. But the analogy is not a claim about
mechanism; it is a claim about formal structure.

The use of Noether's theorem for analyses of conservation in non-physical
systems is another example. Noether's theorem, in its physical form,
asserts that every continuous symmetry of the Lagrangian corresponds
to a conserved quantity. Applied to physical systems, it gives
conservation of energy (from time-translation symmetry), momentum
(from spatial translation), and angular momentum (from rotation).
Applied formally to non-physical systems --- economic systems, social
systems --- it can give conservation laws that are structurally
identical to the physical ones without claiming that the underlying
mechanism is mechanical.

\subsection{What Level 2 commits us to}

A Level 2 treatment commits us to the structural claim and to its
formal consequences, but not to the underlying physical mechanism.

Concretely: if we treat the claim ``rizq is conserved over a
lifetime'' at Level 2, by formalizing it as $\int_0^T R(t)\,dt = K$
where $R(t)$ is the instantaneous provision rate and $K$ is a
fixed constant, then we are committed to the following propositions:

\begin{enumerate}[leftmargin=2em]
    \item There exists a quantity $R(t)$ that we are calling
    ``provision rate'' which has integral over the agent's lifetime
    equal to a fixed value $K$.
    \item Variations in $R(t)$ at one time imply compensating
    variations at other times such that the integral is preserved.
    \item This conservation property has formal consequences (e.g.,
    a Noether-type symmetry that we can identify) which can be
    investigated.
\end{enumerate}

We are not committed to:

\begin{enumerate}[leftmargin=2em]
    \item The proposition that provision is governed by a Lagrangian
    in the physicist's sense.
    \item The proposition that the conservation arises from a
    physical symmetry.
    \item Any specific causal mechanism that produces the
    conservation.
\end{enumerate}

The Level 2 treatment gives us a rigorous formal claim. It does not
give us a physical theory. The claim can still be empirically
tested --- if longitudinal data exists on individual provision over
lifetimes, we can ask whether the integral is approximately constant
across individuals or whether it varies systematically with effort or
disposition --- but the test is of the formal structure, not of any
specific underlying mechanism.

\subsection{The value of Level 2 treatments}

Why bother with Level 2 if it stops short of mechanism?

Three reasons. First, formal structure is real intellectual content.
Knowing that a phenomenon obeys a conservation law, even without
knowing why, organizes the analysis substantially. The phenomena that
obey conservation laws form a class with shared properties. The
Level 2 treatment puts the phenomenon into the class.

Second, Level 2 treatments often suggest mechanisms. The formal
analogy is a heuristic device for hypothesis generation. The
phenomena that exhibit conservation in physics arise from specific
underlying symmetries; the phenomena that exhibit formal conservation
in economics may arise from analogous symmetries that we have not
yet identified. The formal structure points us toward where to look.

Third, Level 2 treatments are honest about the limits of current
knowledge. To assert a Level 3 treatment when only a Level 2 case
can be made is to over-claim. The taxonomy makes the over-claim
visible. It permits the author to do real work at Level 2 without
pretending the work has done more than it has.

\subsection{Level 2 in this book}

A substantial fraction of the worked treatments in Part III are
Level 2 treatments. The treatment of \arb{barakah} as a field
(\xref{ch:two-failures} previewed; full treatment in Chapter 15) is
Level 2: we adopt the formal vocabulary of field theory to organize
the analysis without claiming that barakah is literally an
electromagnetic field. The treatment of \arb{rizq} as a
conservation law (Chapter 10) is Level 2: we adopt the formal
vocabulary of Noether's theorem without claiming that rizq is
literally a Noetherian charge. The treatment of \arb{tajalli} as
continuous creation in the formalism of dynamical systems (Chapter 9)
is Level 2.

The Level 2 treatments are valuable. They organize phenomena that
have not been organized before. They suggest hypotheses. They
generate empirical predictions that the formal analogy entails. They
are not, however, mechanism claims. The book's treatments at Level 2
are flagged as such throughout.

\section{Level 3: Causal mechanism}
\label{sec:level-3}

\begin{definition}[Level 3: Causal mechanism]
\label{def:causal-mechanism}
A metaphysical claim is treated at \emph{Level 3} when the claim is
shown to follow from an actual causal mechanism whose operation can
be specified, derived from first principles or established empirical
regularities, and whose consequences are mathematically demonstrable
and empirically testable. The treatment does not merely import
formal structure from an analogue system; it identifies the actual
process by which the claim is true.
\end{definition}

Level 3 is the strongest claim a treatment can make. It asserts that
the metaphysical statement is a description of a real mechanism, and
it specifies what that mechanism is.

The exemplar of a Level 3 treatment in this book is the treatment of
\arb{riba} (interest) as financial entropy. The Quranic claim ---
that Allah destroys riba and increases charities (\Quran{2:276}) ---
is a metaphysical claim about the dynamics of wealth. A Level 1
treatment would render it as moral exhortation: do not engage in
interest because Allah will not bless your wealth. A Level 2
treatment would render it as a formal analogy: the structure of
interest-based dynamics is structurally similar to the structure of
entropy in thermodynamics. A Level 3 treatment goes further. It
identifies the actual mechanism by which interest destroys real
wealth: the divergence of nominal from real wealth under
multiplicative compound dynamics, the absorbing-barrier dynamics that
make individual experience differ from ensemble expectation, the
non-linear collapse dynamics that follow when the divergence exceeds
sustainable bounds.

This treatment, which Chapter 13 develops in detail, is Level 3
because:

\begin{enumerate}[leftmargin=2em]
    \item It identifies a specific mechanism: compound nominal
    growth at rate $r$ outpacing real growth bounded by $g < r$,
    producing a gap $G(t) = N(t) - P(t) \to \infty$.
    \item The mechanism is derivable from first principles of
    standard finance and macroeconomics. The mathematics is
    well-established (Fisher 1933, Minsky 1986, Keen 2013).
    \item The mechanism produces empirically testable predictions
    (debt-to-GDP ratios above specific thresholds predict financial
    crises; this prediction has been extensively tested).
    \item The Quranic claim ``Allah destroys riba'' is then
    interpretable as an accurate description of this mechanism, not
    as an arbitrary moral rule, not as a metaphor, and not as a
    formal analogy. The mechanism is real and the Quranic statement
    correctly identifies it.
\end{enumerate}

The Level 3 treatment is the strongest possible reading of the
metaphysical claim. The reading is defensible because the mechanism
is well-established in the secular literature and because the match
between the mechanism and the metaphysical statement is exact.

\subsection{What Level 3 commits us to}

A Level 3 treatment commits us to:

\begin{enumerate}[leftmargin=2em]
    \item The existence of the specific mechanism identified.
    \item The truth of the empirical predictions the mechanism
    produces.
    \item The accuracy of the metaphysical statement as a
    description of the mechanism.
\end{enumerate}

The first commitment is empirical: the mechanism either operates or
it does not, and we have specified what it would look like for it to
operate. The second is also empirical: the predictions either bear
out or they do not. The third is interpretive: it is the claim that
the metaphysical statement, when properly understood, is an accurate
description of what the mechanism does.

The third commitment is the most intellectually delicate. It is the
claim that the believer who reads ``Allah destroys riba'' and the
economist who derives the Minsky collapse condition are looking at
the same phenomenon under different descriptions. This is a strong
claim. It is not entailed by either the metaphysical statement alone
or the formal derivation alone; it requires showing that the two
descriptions are coherent and that they refer to the same thing.

The book makes this case carefully and explicitly when claiming
Level 3 treatments. The reader is not asked to accept the
identification on the author's authority; the case is made.

\subsection{Level 3 in this book}

Several treatments in Part III are Level 3, in the author's
assessment, on the basis of the work that has been done up to the
present time. These include:

\begin{itemize}[leftmargin=2em]
    \item The treatment of \arb{riba} as financial entropy
    (Chapter 13).
    \item The treatment of \arb{halal} earnings versus
    interest-based earnings as ergodic versus non-ergodic dynamics
    (Chapter 12).
    \item The treatment of \arb{shukr} (gratitude) and
    its multiplication of wealth as a behavioral mediation chain
    (Chapter 14).
    \item The treatment of the four anti-concentration mechanisms
    (zakat, inheritance partition, riba prohibition, waqf) as a
    coordinated package that stabilizes long-run wealth distribution
    (Chapter 17).
\end{itemize}

The other treatments in Part III are Level 2, with reasons stated in
each chapter.

\section{Worked examples: calibrating intuition}
\label{sec:worked-examples}

The taxonomy is best understood by working through examples. The
following five cases are chosen to span the range of treatments that
the rest of the book will employ.

\subsection{Example 1: the second law of thermodynamics}

Consider the second law of thermodynamics, which asserts that the
entropy of an isolated system never decreases.

Statement: in any closed system, the entropy $S$ satisfies
$dS/dt \geq 0$.

Treatment: Level 3 in physics. The mechanism is statistical: the
increase of entropy reflects the overwhelming numerical preponderance
of high-entropy microstates over low-entropy ones. Boltzmann's
$H$-theorem provides a derivation under specific assumptions. The
mechanism is empirical: it makes predictions that have been verified
across the entire range of physical systems.

Treatment: Level 2 when applied formally to non-physical systems.
The phrase ``social entropy'' or ``organizational entropy,'' applied
to social systems, is a formal analogy. The phenomena being labeled
do not literally have entropy in the Boltzmannian sense; the formal
structure of the claim is borrowed.

Treatment: Level 1 when used as metaphor. ``All things tend to
disorder'' as a general life observation is a Level 1 use of the
second law's vocabulary. No formal claim is being made; no derivation
is being performed.

The second law thus admits all three levels of treatment depending
on context.

\subsection{Example 2: the prohibition of riba}

Consider the Quranic prohibition of riba.

Treatment: Level 1 if presented as moral rule. ``Do not charge
interest because Allah forbids it.'' This is the homiletic register;
no formal claim is made about why the prohibition exists or what it
implies for the dynamics of wealth.

Treatment: Level 2 if presented as formal analogy to known dynamics.
``Interest-based wealth dynamics resemble entropy in thermodynamics
in the following respects.'' The structural claim is rigorous; the
underlying mechanism is not asserted.

Treatment: Level 3 if presented as identification of a specific
financial mechanism. ``The compound growth of nominal claims at rate
$r > g$ (real growth) produces a divergent gap $G(t)$ that, by the
Minsky collapse theorem, requires non-linear correction; the
prohibition of riba is the legal-ethical expression of the
recognition that this mechanism, left unconstrained, destroys real
wealth.'' Here the mechanism is identified, derived, and empirically
testable.

The book treats riba at Level 3 in Chapter 13.

\subsection{Example 3: the children-bring-rizq claim}

Consider the claim ``every child is born with their own rizq''
(\Quran{17:31}).

Treatment: Level 1 if presented as consolation. ``Do not fear
poverty when you have many children; Allah will provide.'' Pastoral;
no theoretical claim.

Treatment: Level 2 if presented as a formal analogy to network
scaling. ``Households with many members have wealth dynamics
structurally similar to other super-linear social systems.'' The
formal structure is suggested; the specific mechanism is not.

Treatment: Level 3 if presented as identification of a specific
network mechanism. ``Households satisfying conditions
$\mathbb{I}_1, \dots, \mathbb{I}_5$ exhibit super-linear wealth
scaling with exponent $\beta > 1$ for the reasons that complex-systems
theory has identified for general super-linear social systems; the
addition of a member therefore increases per-capita wealth, with the
quantitative magnitude derivable from the scaling exponent.''

The book treats this at Level 3 in Chapter 11. The treatment was
sketched in the foundational outline; the full derivation appears in
the relevant chapter when written.

\subsection{Example 4: the qadar/ikhtiyar tension}

Consider the classical theological tension between divine decree
(qadar) and human choice (ikhtiyar).

Treatment: Level 1 in pastoral exhortation. ``Trust in Allah's
decree but tie your camel.'' Wise advice; no theoretical resolution
of the tension.

Treatment: Level 2 by appeal to deterministic interpretations of
quantum mechanics. ``The Bohmian formulation, in which particles have
definite trajectories guided by a non-local pilot wave, dissolves the
apparent contradiction between determinism and the reality of
motion.'' The structural insight is rigorous; the claim that human
choice is governed by Bohmian-type dynamics is not asserted.

Treatment: Level 3 attempted but not (in the author's view)
achievable. The claim that the mind operates by the same physical
mechanism as the Bohmian pilot wave is not establishable on present
knowledge. The Level 2 treatment removes the philosophical
contradiction but does not establish the mechanism.

The book treats qadar at Level 2 in Chapter 16, with explicit
discussion of why Level 3 is not currently available.

\subsection{Example 5: the claim that gratitude multiplies wealth}

Consider the Quranic statement that gratitude is conditional for
increase (\Quran{14:7}).

Treatment: Level 1 in homiletic context. ``Be grateful and Allah
will bless you.''

Treatment: Level 2 by appeal to general well-being literature.
``Grateful people have better outcomes; this resembles the kind of
relationship the Quran describes.'' Suggestive but not derived.

Treatment: Level 3 by specification of the behavioral mediation
chain. ``Gratitude practice produces measurable changes in cortisol
and stress response, which produce measurable changes in
decision-making quality and time horizons; these changes produce
measurable changes in savings rates, investment behavior, and social
capital; the cumulative effect is a measurable increase in long-run
wealth outcomes; each link in the chain is empirically supported.''

The book treats this at Level 3 in Chapter 14. The treatment is
careful: each link in the chain has its own empirical support, and
the chain as a whole produces a cumulative effect that the Quranic
statement is correctly identifying.

\section{Falsifiability and what it requires at each level}
\label{sec:falsifiability}

Karl Popper's criterion of demarcation --- that a theory is
scientific only if it makes predictions that could in principle be
falsified --- has been substantially refined since its 1934
introduction\footnote{Popper, K. (1959), \emph{The Logic of Scientific
Discovery}, Hutchinson; revised editions and subsequent work by
Lakatos, Feyerabend, Kuhn, and others have qualified the criterion
significantly. The author treats falsifiability as a useful but
non-absolute test of theoretical seriousness, in line with the
Lakatosian view of research programs.} but it remains a useful test
of theoretical seriousness. The taxonomy of levels can be
characterized in terms of the falsifiability requirements at each
level.

\textbf{Level 1 treatments are not directly falsifiable.} A Level 1
treatment makes no precise empirical claim. The simile ``rizq is
like air'' is neither true nor false in the empirical sense; it is
suggestive or unsuggestive, evocative or hollow, but it does not
admit of refutation. This is why Level 1 treatments belong to the
homiletic register: they do work that does not require falsifiability
to do.

\textbf{Level 2 treatments are partially falsifiable.} A Level 2
treatment makes a structural claim --- the phenomenon under analysis
shares the formal structure of a known system --- which has empirical
content. The empirical content is the prediction that the phenomenon
will exhibit the relations among variables that the formal structure
implies. If the phenomenon does not exhibit these relations, the
Level 2 treatment is refuted as a structural claim, even if it
remains valid as a metaphor. The Level 2 treatment of barakah as a
field, for example, predicts that the welfare residual after material
controls will exhibit certain structural properties (extension in
space, modulation by intervening factors, etc.). If the welfare
residual is purely random, the Level 2 treatment is falsified at
the structural level.

\textbf{Level 3 treatments are fully falsifiable.} A Level 3
treatment makes a mechanism claim with specific empirical predictions.
The treatment is falsified if the empirical predictions fail. The
treatment of riba as financial entropy predicts specific
relationships between debt-to-GDP ratios and the probability of
financial crisis; these relationships have been extensively tested
and the predictions hold robustly. If the predictions had failed,
the Level 3 treatment would have been refuted.

This pattern --- increasing falsifiability with increasing level ---
is characteristic of the taxonomy and is a virtue of the framework.
A claim treated at Level 3 is the strongest claim and accordingly
exposes itself to the most serious empirical risk. A claim treated at
Level 2 is weaker but still exposed. A claim treated at Level 1 is
the weakest and is largely insulated from refutation.

\subsection{What honest engagement requires}

The taxonomy implies a discipline. The discipline is: \emph{always
treat a claim at the highest level the evidence supports, but never
higher}. To treat a claim at Level 3 when only Level 2 is supported
is to over-claim --- to assert a mechanism without having
established one. To treat a claim at Level 1 when Level 3 is
supported is to under-claim --- to refuse the rigorous treatment that
would be possible.

Honest engagement requires the author to:

\begin{enumerate}[leftmargin=2em]
    \item Specify the level of treatment being attempted.
    \item Provide the appropriate evidence and derivation for that
    level.
    \item Acknowledge when the level falls short of what one might
    have hoped.
    \item Resist the temptation to inflate Level 2 treatments into
    Level 3 claims by rhetorical means.
\end{enumerate}

The book attempts this discipline throughout. The reader who
encounters a Level 3 claim is entitled to demand the full mechanism
derivation; the reader who encounters a Level 2 claim should
recognize that the formal structure is being borrowed without
mechanism commitment; the reader who encounters a Level 1 expression
should treat it as motivation, not as substantive treatment.

\section{The special role of Akbarian metaphysics}
\label{sec:akbarian-role}

A note on the special status of the Akbarian tradition in the
taxonomy.

The Akbarian metaphysics --- the system articulated by Ibn Arabi and
developed by Mulla Sadra, Shah Wali Allah, and others --- does not
fit neatly into any of the three levels as defined above. It is not
metaphor (Level 1): it is articulated with extraordinary precision,
it derives propositions from premises, and it is responsive to
counter-arguments. It is not formal analogy (Level 2): it does not
import the formal structure of a different system, but rather
develops its own internal formal apparatus. It is not causal
mechanism in the modern empirical sense (Level 3): the mechanisms it
describes are not the kind that contemporary science recognizes as
candidates for empirical test.

The Akbarian tradition operates at a level that the three-level
taxonomy does not directly capture. The author proposes to call this
\emph{ontological articulation}: the precise specification of the
structure of being itself, expressed in vocabulary that does not
reduce to either the formal sciences or the empirical sciences as
currently constituted.

The relationship between ontological articulation and the three
levels is the following. The three levels are levels of
\emph{theoretical} engagement with phenomena, taking the underlying
ontology as given. Ontological articulation is the prior work of
specifying what kind of thing the phenomena are. The Akbarian
tradition, by articulating the ontology of \arb{wujud} and its
\arb{tajalli}, provides the ontological substrate that the
theoretical work of Levels 2 and 3 then operates on.

This is not unique to Islamic metaphysics. Every theoretical
discipline rests on an ontology that the discipline does not itself
articulate. Modern physics rests on a materialist or naturalist
ontology that physics qua physics does not derive; it is presupposed.
The genuine ontological work in modern science is done by the
philosophy of science, the philosophy of physics, the philosophy of
mind. The same is true here. The Akbarian tradition does the
ontological work that the theoretical engagement with the
metaphysical claims then operates on.

The book's chapter on the Akbarian tradition (Chapter 9) accordingly
treats it as a system of ontological articulation rather than as a
candidate for the three-level taxonomy. The Akbarian framework is
not classified as Level 1 or 2 or 3; it is the ontological substrate
that the three levels operate on. Understanding this is important
for understanding why Chapter 9 has the structure it does and why
the chapters that follow it are able to speak in the way they do.

\section{Looking ahead: how the levels apply in Part III}
\label{sec:ahead-2}

The eight worked treatments in Part III each claim a specific level
of treatment. The reader who keeps the taxonomy in mind will find
the chapters easier to evaluate. A preview:

\begin{itemize}[leftmargin=2em]
    \item \textbf{Chapter 10. Rizq as conserved quantity}: Level 2
    primarily, with some Level 3 elements where the empirical
    correlates can be specified. The conservation of rizq as an
    integral over the lifetime is a structural claim with formal
    consequences; the specific mechanism by which the conservation
    is enforced is not asserted.
    \item \textbf{Chapter 11. The network of provision}: Level 3.
    The specific mechanism is super-linear network scaling under
    Islamic structural conditions; the mechanism is established in
    the network-economics literature; the empirical predictions are
    testable with available household data.
    \item \textbf{Chapter 12. Halal versus haram dynamics}: Level 3.
    The mechanism is ergodicity (or lack thereof) of multiplicative
    wealth processes with absorbing barriers; the mechanism is
    established in ergodicity economics; the empirical predictions
    follow rigorously.
    \item \textbf{Chapter 13. Riba as financial entropy}: Level 3.
    The mechanism is the Fisher-Minsky-Keen dynamics of debt-driven
    economies; the mechanism is well-established; the empirical
    predictions are extensively tested.
    \item \textbf{Chapter 14. Gratitude and the multiplication of
    wealth}: Level 3. The mechanism is a four-link behavioral
    mediation chain (cortisol, decision quality, social capital,
    savings rate); each link is established in the behavioral
    economics literature; the cumulative effect is testable.
    \item \textbf{Chapter 15. Barakah as operational field}: Level 2
    primarily, with a research program for moving toward Level 3.
    The structural claim --- that the welfare residual after material
    controls has systematic content --- is operational; the specific
    mechanisms by which the systematic content arises are subject to
    empirical investigation.
    \item \textbf{Chapter 16. Qadar and the resolution of
    determinism}: Level 2. The Bohmian and QBist formulations
    dissolve the apparent contradiction between qadar and ikhtiyar
    by formal structure; the claim that the underlying physics is
    Bohmian (rather than, say, Many-Worlds) is not asserted.
    \item \textbf{Chapter 17. Architecture of permanence}: Level 3.
    The mechanism is the dynamics of preferential attachment as
    modulated by the four Islamic anti-concentration institutions;
    the mechanism is well-established in the wealth-distribution
    literature; the empirical predictions follow from agent-based
    simulation of the full package.
\end{itemize}

Of the eight, six are Level 3 and two are Level 2. The Level 2
treatments (rizq conservation, qadar) are not weaknesses of the
project; they are honest acknowledgments of what current knowledge
permits. Future research may move them to Level 3.

The classification will be important when we evaluate the cumulative
case the book makes. A book composed entirely of Level 1 treatments
would be a homiletic work. A book composed entirely of Level 2
treatments would be a structural analysis. A book that mixes Level 2
and Level 3 treatments, with the level of each treatment explicitly
flagged, is what the present work attempts.

\begin{summarybox}[Chapter summary]
\textbf{The three-level taxonomy of rigor:}
\begin{itemize}[leftmargin=1.5em]
    \item \textbf{Level 1: Metaphor.} Figurative or pedagogical
    illustration. Not directly falsifiable. Appropriate to the
    homiletic register.
    \item \textbf{Level 2: Formal analogy.} Mathematical structure
    borrowed from a known system without claiming the underlying
    mechanism. Partially falsifiable (the formal structure either
    holds or does not).
    \item \textbf{Level 3: Causal mechanism.} Identification of an
    actual mechanism, derived from first principles or established
    empirical regularities. Fully falsifiable.
\end{itemize}

\textbf{The discipline}: always treat a claim at the highest level
the evidence supports, but never higher.

\textbf{The Akbarian tradition} is best understood not as one of the
three levels but as \emph{ontological articulation} that provides
the substrate on which the levels operate.

\textbf{In Part III}: six of the eight worked treatments aim at
Level 3, two at Level 2; each chapter explicitly flags its level.

\textbf{What comes next}: Chapter 3 surveys the textual foundations
on which all of the worked treatments will draw.
\end{summarybox}

\cleardoublepage
